\documentclass[11pt]{article}
\usepackage[margin=1in]{geometry}
\usepackage{amsmath,amssymb,amsthm,booktabs,longtable,array}
\usepackage[colorlinks=true,linkcolor=blue,urlcolor=blue]{hyperref}
\usepackage{microtype}
\newtheorem{theorem}{Theorem}[section]
\newtheorem{lemma}[theorem]{Lemma}
\newtheorem{proposition}[theorem]{Proposition}
\newtheorem{corollary}[theorem]{Corollary}
\newcommand{\E}{\mathbb E}
\newcommand{\Var}{\operatorname{Var}}
\newcommand{\Inf}{\operatorname{Inf}}
\newcommand{\sech}{\operatorname{sech}}
\newcommand{\supp}{\operatorname{supp}}
\setlength{\emergencystretch}{2em}
\title{Quadratic polynomial-influence AA for Gibbs laws:\\
centered coefficients, a dimension-free bound,\\and an exact obstruction to the proposed base inequality}
\author{Research continuation in the AAconjecture workspace}
\date{14 November 2026}
\begin{document}
\maketitle

\noindent\textbf{Status.}
This report gives self-contained proofs of three results: the proposed
base inequality (B) is false even for zero-field independent Ising pairs;
one-sided uncentred marginal-product $L^2$ domination already fails in
degree two; and a different, centered and variance-weighted coefficient
argument proves polynomial-influence AA in degree at most two, with
arbitrary finite fields, interaction arities, and coordinate supports,
when the weighted interaction strength is at most a fixed constant below
one. A pairwise corollary covers the full coupling-based Dobrushin region
with a fixed margin. The all-degree problem, in particular unrestricted
$d\ge3$, is not solved. No independent audit, novelty claim, or Lean
formalization is asserted. Finite checks are supplementary.

The requested prompt and handoff were read from \texttt{prompt/}.
Their referenced \texttt{.ultramath/} archive is absent in this workspace,
so its original manuscripts and verification scripts could not be read
or run. Accordingly none of their claimed theorems is used without a
proof here. The supplied \texttt{gibbs\_history\_aa\_corrected.pdf}
is used to preserve the distinction between query simulation and
physical polynomial influence; its query theorem is not an input to
the positive theorem below. Existing source proofs are left unchanged.

\section{Statement, quantifiers, and conventions}

Let $n$ be arbitrary and let all coefficients in
\[
 H(x)=-\sum_{|A|\ge2}J_A\chi_A(x)-\sum_i h_ix_i,
 \qquad \mu(x)=Z^{-1}e^{-\beta H(x)},\qquad
 \chi_A(x)=\prod_{i\in A}x_i
\]
be finite, with $\beta\ge0$. Define
\[
 K=\max_i\beta\sum_{A\ni i,\,|A|\ge2}|J_A|,
 \qquad
 \omega=\max_i\beta\sum_{A\ni i,\,|A|\ge2}(|A|-1)|J_A|.
\]
Thus $K\le\omega$. The field parameter
$\Lambda_{\rm fld}=\max_i\beta|h_i|$ is unrestricted and may grow
with $n$. No support, hitting-set, degree-of-graph, or arity restriction
is imposed.

For a coordinate $i$, put
\[
 m_i=\E_\mu X_i,\quad \sigma_i^2=1-m_i^2,\quad
 \theta_i(x_{-i})=\mu(X_i=1\mid x_{-i}),\quad
 v_i(x_{-i})=4\theta_i(1-\theta_i).
\]
Here $m_i$ denotes a marginal mean, not a conditional mean. All
$\sigma_i>0$ because the Gibbs law has full support. For
$\partial_i p=(p(+1,x_{-i})-p(-1,x_{-i}))/2$, the conventions are
\begin{align*}
 \Inf_i^{\rm fl}(p)&=4\E(\partial_i p)^2,\\
 \Inf_i^{\rm pv}(p)&=\E[v_i(\partial_i p)^2],\\
 \Inf_i^{\rm rs2}(p)&=2\Inf_i^{\rm pv}(p).
\end{align*}
Consequently $\Inf_i^{\rm fl}\ge4\Inf_i^{\rm pv}
=2\Inf_i^{\rm rs2}$. No TV/pivotal influence is used below.

Define the TV Dobrushin matrix by $c_{ii}=0$ and
\[
 c_{ij}=\sup_{x,z:\,x_{-j}=z_{-j}}
       |\theta_i(x_{-i})-\theta_i(z_{-i})|\quad(i\ne j),
 \qquad r=\max_i\sum_j c_{ij}.
\]
This uses differences of probabilities, not twice those differences
as in the conditional-mean convention of parts of the handoff.

\begin{theorem}[Quadratic AA with arbitrary fields]
\label{thm:main}
Fix $K_0<\infty$, $0\le r_0<1$, and $0<v\le1$ before choosing
$n$. For every Gibbs law above with $K\le K_0$ and $r\le r_0$,
every real multilinear polynomial of degree at most two satisfying
$\|p\|_{\infty,\{-1,1\}^n}\le1$ and $\Var_\mu p\ge v$ has
a coordinate $i$ for which simultaneously
\begin{align}
 \Inf_i^{\rm fl}(p)&\ge
       \frac{e^{-16K_0}(1-r_0)^6}{93312}\,v^3,
       \label{eq:main-fl}\\
 \Inf_i^{\rm pv}(p)&\ge
       \frac{e^{-20K_0}(1-r_0)^6}{373248}\,v^3,
       \label{eq:main-pv}\\
 \Inf_i^{\rm rs2}(p)&\ge
       \frac{e^{-20K_0}(1-r_0)^6}{186624}\,v^3.
       \label{eq:main-rs}
\end{align}
All constants are independent of $n$, field strengths, and support size.
The flip estimate is obtained directly and is slightly stronger than
four times the displayed posterior-variance estimate.
\end{theorem}

\begin{corollary}[The requested weighted-strength regime]
\label{cor:omega}
For every fixed $0\le\omega_0<1$, the preceding theorem holds
uniformly for $\omega\le\omega_0$, with $K_0=r_0=\omega_0$.
In particular, at $\omega\le1/8$,
\[
 \boxed{\max_i\Inf_i^{\rm fl}(p)
       \ge \frac{e^{-2}(7/8)^6}{93312}\,v^3.}
\]
This statement allows odd as well as even interaction arities and
arbitrary finite fields. It remains true under arbitrary pinnings,
after removing pinned coordinates.
\end{corollary}

\begin{corollary}[Pairwise Dobrushin region]
\label{cor:dobrushin}
For pairwise Ising Hamiltonians with arbitrary fields and signs,
assume, uniformly in $n$,
\[
 \alpha:=\max_i\sum_{j\ne i}\tanh(\beta|J_{ij}|)
       \le1-\kappa,\qquad 0<\kappa\le1.
\]
Then for every globally bounded real degree-at-most-two polynomial
of variance at least $v$,
\[
 \max_i\Inf_i^{\rm fl}(p)
 \ge\frac{\kappa^{14}}{93312(2-\kappa)^8}\,v^3.
\]
The posterior-variance lower bound is
$\kappa^{16}v^3/[373248(2-\kappa)^{10}]$.
The conclusion is uniform under pinnings. This is an influence theorem
for quadratics, not a claim about every polynomial degree or every
quantum acceptance function in the Dobrushin region.
\end{corollary}

\section{The base inequality fails: exact pair-block calculations}

The default question in the prompt asks whether
\[
 \widehat q(\emptyset)^2\le\lambda_0^{-1}\E_\mu q^2
 \tag{B}
\]
can hold with a fixed $\lambda_0>0$ for all bounded quadratics in
zero-field high-temperature Gibbs families. The coefficient on the
left is $\E_U q$, where $U$ is uniform, not $\E_\mu q$.

Take $n=2m$, $\beta=1$, and
\[
 H_m(x)=-a\sum_{j=1}^m x_{2j-1}x_{2j},\qquad
 a=\tfrac12\log(5/4)<\tfrac18.
\]
There are no fields, every coordinate has one incident pair term,
and $K=\omega=a$. Put $Y_j=X_{2j-1}X_{2j}$ and
$\rho=\tanh a=1/9$. The $Y_j$ are independent with mean $\rho$,
while every physical one-site marginal is uniform. Indeed each pair
has mass $(1+\rho xy)/4$ and partition function $4\cosh a$;
the full partition function is $(4\cosh a)^m$.

The globally bounded quadratic
\[
 q_m(x)=\frac{m^{-1}\sum_jY_j-\rho}{1+\rho}
       =\frac9{10m}\sum_jY_j-\frac1{10}
\]
has exact statistics
\begin{align}
 \|q_m\|_\infty&=1,&\widehat q_m(\emptyset)&=-1/10,
       &\E_\mu q_m&=0,\nonumber\\
 \E_\mu q_m^2=\Var_\mu q_m&=\frac4{5m},&
 \E_Uq_m^2&=\frac1{100}+\frac{81}{100m}.\label{eq:pair-stats}
\end{align}
Thus (B) would require $\lambda_0\le80/m$ for every $m$,
which is impossible. This construction works at arbitrarily small
fixed positive interaction strength: for general $0<\rho<1$ the
same ratio is $(1-\rho^2)/(m\rho^2)$.

\begin{proposition}[Sharp degree-two base constant on the pair family]
The best constant in (B), at each fixed $m$ for this law and all
degree-at-most-two polynomials, is
\[
 \lambda_m=\frac{80}{80+m}.
\]
The same optimum holds if one requires $\|q\|_\infty\le1$.
\end{proposition}
\begin{proof}
Write $Z_j=Y_j-\rho$. The degree-two space has the decomposition
\[
 q=c+\sum_j b_jZ_j+R,
\]
where $R$ is a linear combination of single spins and products of
spins from distinct pairs. Under the pair law, $R$ is orthogonal
to $1,Z_1,\ldots,Z_m$: flipping both spins of any involved pair
changes the relevant odd factor and preserves the measure.
Also $\E_U R=0$ and $\E_U Z_j=-\rho$. Hence
\[
 (\E_U q)^2=(c-\rho\sum_jb_j)^2
 \le\left(1+\frac{m\rho^2}{1-\rho^2}\right)
       \left(c^2+(1-\rho^2)\sum_jb_j^2\right)
 \le\left(1+\frac{m\rho^2}{1-\rho^2}\right)\E_\mu q^2.
\]
Equality is attained, after harmless scalar normalization, by
\[
 q_m^*(x)=\frac{80+m-9\sum_jY_j}{80+10m},\qquad
 \|q_m^*\|_\infty=1.
\]
Its uniform coefficient is $(80+m)/(80+10m)$, its Gibbs mean is
$80/(80+10m)$, its variance is $80m/(80+10m)^2$, and its squared
$L^2(\mu)$ norm is $80(80+m)/(80+10m)^2$. Their ratio proves
sharpness. Homogeneity explains why boundedness does not change
the optimal constant.
\end{proof}

These are not AA counterexamples. For $q_m$ in
\eqref{eq:pair-stats}, every coordinate has
\[
 \Inf_i^{\rm fl}=\frac{81}{25m^2},\quad
 \Inf_i^{\rm pv}=\frac4{5m^2},\quad
 \Inf_i^{\rm rs2}=\frac8{5m^2}.
\]
In particular $\Inf_i^{\rm fl}=(81/16)(\Var_\mu q_m)^2$,
and the variance tends to zero. Fixed-$v$ quantifiers cannot be
violated by this sequence.

The same family resolves the proposed uncentred one-sided $(H_d)$
input negatively already at $d=2$. For $f_m=m^{-1}\sum_jY_j$,
the marginal product law is $\nu=U$ and
\[
 \|f_m\|_\infty=1,\qquad
 \frac{\E_\mu f_m^2}{\E_\nu f_m^2}
 =\frac{m+80}{81}\longrightarrow\infty.
\]
No $n$-uniform $A_2$ exists for the stated uncentred domination.
The positive proof below neither assumes (B) nor transfers influences
through this false moment comparison.

\section{Self-contained analytic inputs}

\subsection{Dobrushin contraction supplies the needed Poincare inequality}

Let $P_i f=\E[f\mid X_{-i}]$, regarded as a function on the full
cube, and $P=n^{-1}\sum_iP_i$. Each $P_i$ is an orthogonal
projection in $L^2(\mu)$; $P$ is self-adjoint and positive
semidefinite. Write
$\delta_i(f)=\max_x|f(x)-f(x^{(i)})|$.
For $i\ne j$, a maximal coupling of the two conditional laws of
spin $j$ at inputs differing only in $i$ gives
\[
 \delta_i(P_j f)\le\delta_i(f)+c_{ji}\delta_j(f),
 \qquad \delta_i(P_i f)=0.
\]
Indeed, the coupled configurations always differ in $i$, and
additionally differ in $j$ with probability at most $c_{ji}$.
The triangle inequality bounds their function difference by the
displayed oscillations. Consequently
\[
 \sum_i\delta_i(Pf)\le
 \left(1-\frac{1-r}{n}\right)\sum_i\delta_i(f).
\]
Every nonconstant eigenfunction of $P$ has positive oscillation
seminorm, so its eigenvalue is at most $1-(1-r)/n$.
Diagonalizing this finite self-adjoint operator proves
\begin{equation}
 \Var_\mu f\le C\sum_i\E\Var(f\mid X_{-i}),
 \qquad C=\frac1{1-r}.
 \label{eq:poincare}
\end{equation}
The Dirichlet form identity used here is
$\langle f,(I-P)f\rangle=n^{-1}\sum_i\E\Var(f\mid X_{-i})$.
This is not an application of unordered independent Efron--Stein
to a dependent measure; the factor $C$ has just been proved.

For later use, applying \eqref{eq:poincare} to a linear function
and using total variance separately at each coordinate gives
\begin{equation}
 \Var_\mu\left(\sum_jt_jX_j\right)
 \le C\sum_jt_j^2\E v_j
 \le C\sum_jt_j^2\sigma_j^2.
 \label{eq:weighted-cov}
\end{equation}
There is no assertion that the spins are independent.

For the Gibbs specification, changing $x_j$ changes the local
field at $i$ by at most $2\beta\sum_{A\ni i,j}|J_A|$.
The derivative of $(1+\tanh z)/2$ is at most $1/2$. Therefore
$c_{ij}\le\beta\sum_{A\ni i,j}|J_A|$, and $r\le\omega$.
This proves the contraction hypothesis in Corollary~\ref{cor:omega}
with the stated TV normalization.

\subsection{Conditional and marginal variances are comparable without a field bound}

The conditional log odds at site $i$ are
\[
 \ell_i(x_{-i})=2\beta h_i+
              2\beta\sum_{A\ni i,\,|A|\ge2}J_A\chi_{A\setminus\{i\}}(x),
\]
and their oscillation is at most $4K$. The marginal probability
$\bar\theta_i=\E\theta_i$ lies between the minimum and maximum
conditional probabilities; its log odds lie in the same interval.
The function $g(\ell)=4e^\ell/(1+e^\ell)^2$ satisfies
$|(\log g)'(\ell)|\le1$. Hence, with $D=e^{4K}$,
\begin{equation}
 D^{-1}\sigma_i^2\le v_i(x_{-i})\le D\sigma_i^2
 \quad\hbox{for every }i,x_{-i}.
 \label{eq:variance-comparison}
\end{equation}
Large fields may make both sides very small, but introduce no
field-dependent constant. This is a local conditional comparison,
not a global density-ratio bound or an atom floor.

\subsection{An elementary quadratic coefficient bound}

Write
\[
 p(x)=a_0+\sum_i a_ix_i+\sum_{i<j}b_{ij}x_ix_j,
 \qquad b_{ji}=b_{ij},\quad b_{ii}=0.
\]
Its odd part is $\sum_i a_ix_i$, of supremum norm at most one,
so $\sum_i|a_i|\le1$. Its quadratic homogeneous part $Q$ has
supremum norm at most two: the even part has norm at most one and
$|a_0|=|\E_U p|\le1$.

Let $A_{ij}=b_{ij}/2$ for $i\ne j$ and $A_{ii}=0$, and let
$B(x,y)=x^TAy$. Since a multilinear polynomial on $[-1,1]^n$
attains its extrema at vertices, polarization gives
\[
 B(x,y)=Q((x+y)/2)-Q((x-y)/2),\qquad
 \|B\|_{[-1,1]^n\times[-1,1]^n}\le4.
\]
The following elementary proof avoids importing a quantitative
Bohnenblust--Hille theorem. Put
$R_i=(\sum_j|A_{ij}|^2)^{1/2}$ and
$S_j=(\sum_i|A_{ij}|^2)^{1/2}$. Three-factor Holder gives
\[
 \sum_{ij}|A_{ij}|^{4/3}
 \le\left(\sum_iR_i\right)^{2/3}
       \left(\sum_jS_j\right)^{2/3}.
\]
To verify it, factor each summand as
$(|A_{ij}|^2/R_i)^{1/3}(|A_{ij}|^2/S_j)^{1/3}
(R_iS_j)^{1/3}$, omitting zero rows and columns.
For independent uniform signs $\varepsilon_j$ and
$Z=\sum_jd_j\varepsilon_j$, one has
$\E Z^2=\sum_jd_j^2$ and $\E Z^4\le3(\sum_jd_j^2)^2$.
Holder yields
$\E|Z|\ge(\sum_jd_j^2)^{1/2}/\sqrt3$.
Applying this rowwise and maximizing over $x$ gives
\[
 \sum_iR_i\le\sqrt3\E_\varepsilon\sum_i
       |\sum_jA_{ij}\varepsilon_j|\le\sqrt3\|B\|\le4\sqrt3,
\]
and the same column bound. Thus
\begin{equation}
 \sum_{i<j}|b_{ij}|^{4/3}
   =2^{1/3}\sum_{ij}|A_{ij}|^{4/3}
   \le 8\,3^{2/3}<17.
 \label{eq:quadratic-coeff}
\end{equation}

Now center at the actual marginal means, without identifying them
with uniform means:
\[
 p=c_0+\sum_i\widetilde a_i(x_i-m_i)
          +\sum_{i<j}b_{ij}(x_i-m_i)(x_j-m_j),\qquad
 \widetilde a_i=a_i+\sum_jb_{ij}m_j.
\]
The constant $c_0$ is not generally $\E_\mu p$.
Because $m\in[-1,1]^n$ and $b=2A$,
\[
 \sum_i|\widetilde a_i|
 \le1+\sum_i|\sum_jb_{ij}m_j|
 \le1+2\max_{x\in\{-1,1\}^n}|B(x,m)|\le9.
\]
Define weighted coefficients
$u_i=\sigma_i\widetilde a_i$ and
$w_{ij}=\sigma_i\sigma_jb_{ij}$. Since $\sigma_i\le1$,
\begin{equation}
 \sum_i|u_i|^{4/3}+\sum_{i<j}|w_{ij}|^{4/3}
 \le9^{4/3}+8\,3^{2/3}<36.
 \label{eq:weighted-coeff}
\end{equation}
The numerical inequality follows already from $9^4<19^3$ and
$8^3 3^2<17^3$. No basis orthogonality under $\mu$ is assumed.

\section{Proof of the quadratic theorem}

The derivative has the exact centered form
\[
 R_i:=\partial_i p=\widetilde a_i+
                      \sum_{j\ne i}b_{ij}(X_j-m_j).
\]
Using \eqref{eq:weighted-cov} gives
\[
 \E R_i^2\le\widetilde a_i^2+C\sum_jb_{ij}^2\sigma_j^2.
\]
Apply \eqref{eq:poincare}, then \eqref{eq:variance-comparison},
and count each quadratic coefficient twice:
\begin{align}
 \Var_\mu p
 &\le C\sum_i\E[v_iR_i^2]
 \le CD\sum_i\sigma_i^2\E R_i^2\nonumber\\
 &\le CD\sum_i u_i^2+2C^2D\sum_{i<j}w_{ij}^2
 \le L\left(\sum_i u_i^2+\sum_{i<j}w_{ij}^2\right),
 \qquad L=2C^2D.\label{eq:weighted-gram}
\end{align}
This is a weighted centered-coefficient estimate, not the false
uncentred comparison $(H_2)$ and not the unweighted Gram assertion
from the handoff.

Let $M=\max\{|u_i|,|w_{ij}|:i<j\}$, including all linear indices.
Equations \eqref{eq:weighted-coeff}--\eqref{eq:weighted-gram} imply
\[
 v\le36LM^{2/3},\qquad M^2\ge\left(\frac{v}{36L}\right)^3.
 \tag{E1}
\]
If $M=|u_i|$, Jensen gives
$\E R_i^2\ge\widetilde a_i^2\ge u_i^2$ and
\[
 \Inf_i^{\rm pv}\ge D^{-1}\sigma_i^2\E R_i^2
       \ge D^{-1}u_i^2\ge D^{-2}M^2,
 \qquad \Inf_i^{\rm fl}\ge4M^2\ge4D^{-1}M^2.
\]
If $M=|w_{ij}|$, condition $R_i$ on all coordinates except $j$.
Its slope in $X_j$ is exactly $b_{ij}$, so
\[
 \E R_i^2\ge\E\Var(R_i\mid X_{-j})
       =b_{ij}^2\E v_j\ge D^{-1}\sigma_j^2b_{ij}^2.
\]
Consequently
\[
 \Inf_i^{\rm pv}\ge D^{-2}w_{ij}^2,
 \qquad
 \Inf_i^{\rm fl}\ge4D^{-1}\sigma_j^2b_{ij}^2
                    \ge4D^{-1}w_{ij}^2.
\]
Thus in both cases the same coordinate satisfies these bounds.
Substitute (E1), $L=2D/(1-r)^2$, and $D=e^{4K}$:
\[
 \Inf_i^{\rm fl}\ge\frac{e^{-16K}(1-r)^6}{93312}v^3,
 \qquad
 \Inf_i^{\rm pv}\ge\frac{e^{-20K}(1-r)^6}{373248}v^3.
\]
Monotonicity in $K,r$ proves Theorem~\ref{thm:main}; the resampling
formula follows by multiplying by two.

For Corollary~\ref{cor:omega}, we already proved $K,r\le\omega$.
Pinning shrinks each interaction support, moves singleton residuals
to the unrestricted fields, and cannot increase $\omega$ after
merging coefficients by the triangle inequality. Thus the same
constant applies after pinnings. Polynomials remain bounded on the
restricted cube and have degree at most two. No rare-event loss occurs.

For Corollary~\ref{cor:dobrushin}, the exact Ising conditional formula
implies $c_{ij}\le\tanh(\beta|J_{ij}|)$: the difference of the
two logistic probabilities is maximized at zero residual field.
Also for $t_j=\tanh(\beta|J_{ij}|)$ with $\sum_jt_j<1$,
\[
 \sum_j\beta|J_{ij}|=\sum_j\operatorname{arctanh}t_j
       \le\operatorname{arctanh}(\sum_jt_j)
       \le\operatorname{arctanh}(1-\kappa).
\]
The first inequality follows termwise from the positive power series
for $\operatorname{arctanh}$ and
$\sum_jt_j^{2\ell+1}\le(\sum_jt_j)^{2\ell+1}$.
Set $r_0=1-\kappa$ and
$K_0=\frac12\log((2-\kappa)/\kappa)$ in the theorem.
Pinning deletes couplings and changes only fields, preserving the
same coupling-based Dobrushin bound.

\section{A necessary scope correction for zero fields and odd arity}

Zero one-site fields alone do not imply $\E X_i=0$ when odd
interactions are allowed. This matters because the handoff's
degree-two unweighted Gram sketch uses that implication.
On four spins, take $\beta=1$ and
\[
 H=-a\chi_{\{1,2,3\}}-a\chi_{\{1,2,3,4\}},\qquad
 a=\tfrac12\log(21/20),\quad t=\tanh a=1/41.
\]
All fields vanish and
$\omega=5a<1/8$, since $\log(21/20)<1/20$.
The exact uniform-density factor is
$(1+t\chi_{123})(1+t\chi_{1234})$, whose expectation is one.
Since $\chi_{123}\chi_{1234}=X_4$,
\[
 \E_\mu X_4=t^2=1/1681\ne0.
\]
Take independent copies of this block and let $Z_1,\ldots,Z_m$
be their fourth spins, with common mean $s=1/1681$. For
$N=\binom m2$ and $p=N^{-1}\sum_{j<k}Z_jZ_k$, one has
$\|p\|_\infty=1$, $\deg p=2$, and
\[
 \frac{\Var_\mu p}{\sum_S\widehat p(S)^2}
 =1-s^4+2(m-2)s^2(1-s^2)\longrightarrow\infty.
\]
Indeed each pair monomial has variance $1-s^4$; two pairs sharing
one index have covariance $s^2(1-s^2)$, and disjoint pairs have
zero covariance. This refutes the unweighted, all-arity version
of the handoff's zero-field degree-two Gram claim. The globally
spin-flip-symmetric case, including zero-field pairwise Ising and
Hamiltonians with only even interactions, is not contradicted.
Theorem~\ref{thm:main} needs no mean-zero hypothesis and covers
this example using \eqref{eq:weighted-gram} instead.

\section{Checks, comparison with previous work, and remaining gaps}

The companion \texttt{verify.py} uses the Python standard library.
Finite Gibbs laws are represented as rational products of factors
$1+t_A\chi_A$ and $1+u_ix_i$, which correspond exactly to
couplings $\operatorname{arctanh}(t_A)$ and fields
$\operatorname{arctanh}(u_i)$ at $\beta=1$.
Moment and influence calculations use exact rational arithmetic;
exponential bounds and fractional coefficient powers use floating
point only in the supplementary inequality comparisons.
It enumerates pair blocks and short cycles, checks the sharp pair
constant and $(H_2)$ obstruction, checks the odd-arity example,
and tests each step of the positive argument on signed interactions,
strong biases, and pinnings. Its actual output is in
\texttt{verification-output.txt}. These tests do not certify a
dimension-uniform theorem; the preceding arguments do.

The pair example is compatible with accepted product-block results:
it refutes a proposed comparison, not AA. Biased products are covered
by the new theorem with $K=r=0$; no uniform marginal atom floor is
required. Finite fields can approach degeneracy uniformly without
changing the bound. Exactly deterministic coordinates may be removed
and the theorem applied to the remaining law when its stated Gibbs
hypotheses hold. No universal assertion for arbitrary zero-mass
interacting laws is made.

The positive theorem settles degree at most two for unrestricted
support in the prompt's $\omega\le1/8$ class. It does not establish
uniform (WA) or CI-avg, recover (B) or $(H_2)$, prove the optimal
variance exponent, or extend the centered-coefficient calculation to
all degrees. At degree three, derivatives are quadratic rather than
affine: controlling their centered moments and extracting weighted
coefficients requires additional work, and the known pair-block
obstructions must still be respected. The general $d\ge3$ target
remains open in this report.

The corrected Gibbs-history PDF proves a different result: its
Theorem 4.2 approximates quantum acceptance probabilities, and its
Section 5 explicitly does not assert polynomial-influence AA.
Its Lemmas 2.2--2.4 distinguish complete-trajectory map independence
from the generally dependent posterior given the final input.
Neither that independence nor the $4T\sqrt\gamma$ hybrid bound
is used here. We obtain physical coordinate influence directly
from derivatives, local conditional variances, and a spectral
contraction estimate.

\appendix
\section{Source and verification ledger}
\begin{enumerate}
\item \texttt{prompt/polynomial-aa-next-agent-prompt.md}: target,
normalizations, default question (B), required scope, and the correction
to the earlier degree-one necessity claim. No necessity claim is used.
\item \texttt{prompt/polynomial-aa-research-handoff.md}: the pair-block
test family, proposed coefficient route, auxiliary moment formulations,
and previously accepted special cases. Its statements are context;
all mathematical inputs used in this report have been re-proved.
\item \texttt{gibbs\_history\_aa\_corrected.pdf}, especially Lemmas
2.2--2.4, Theorem 4.2, and Section 5: scope of query simulation and
the posterior-independence distinction. No theorem is imported from
this PDF into the influence proof.
\item A. Defant, M. Mastylo, A. Perez,
\href{https://arxiv.org/abs/1707.09186}{\emph{Bohr's phenomenon for
functions on the Boolean cube}}, arXiv:1707.09186: primary abstract
consulted for coefficient-inequality background. No constant or theorem
from it is imported; the required degree-two estimate is proved here.
\item The original reviewed-literature file, derivative-route manuscript,
replacement Gram manuscript, and scripts listed under \texttt{.ultramath/}
were unavailable because that directory is absent. The handoff exists
under \texttt{prompt/}, not at the requested \texttt{results/} path.
New verification code was written only after these availability checks
and after reading the supplied materials. No absent script is reported
as executed.
\end{enumerate}
Apart from elementary finite-dimensional spectral theory, Holder,
Cauchy--Schwarz, and conditional expectation identities, there are no
unproved imported mathematical inputs to Theorem~\ref{thm:main}.
This is a written proof with numerical checks, not a proof-assistant
certificate or an independent referee review.
\end{document}
